\section{Data, code, and release status}\label{s:data}

The publication-oriented repository is
\href{https://github.com/Vulkin-prog/vdw-rabung-bounds}{\texttt{Vulkin-prog/vdw-rabung-bounds}}.
It was exported under an explicit allowlist from the development repository at
the immutable source commit
\texttt{02796ae9e08da5b051a94b2d9001763908375f6e}.  At the time of this
revision it remains a private preparation repository, not a citable release.
Its machine-readable status ledger deliberately leaves the certificate,
campaign-recovery, ordered-prime, external-replication, rights, and DOI gates
open.  An immutable archive of the frozen release tag, rather than a moving
branch, will be cited after those gates pass.

The source registry \texttt{audit/claims.json} contains exactly the thirteen
succinct Rabung certificates $(p,r,k)$.  The separate
\texttt{audit/baselines.json} records published, archived, and current strict
lower bounds with source locators.  Running
\texttt{tools/close\_baselines.py} produces the recurrence-closed bounds,
their provenance DAG, and the \LaTeX{} table used in
Section~\ref{s:bounds}.  This prevents a manually maintained record table from
being compared with an unclosed bibliography.

For every direct claim, the final layout is
\[
 \texttt{results/claims/<claim-id>/\{manifest.json,raw/*\}}.
\]
The manifest records the claim, verdict, max-run and (B$^*$) profiles, Git
commit, source and binary hashes, build commands, compiler versions, and the
SHA-256 and byte length of every raw standard stream.  The fail-closed parser
in \texttt{tools/claim\_audit.py} validates those manifests and generates the
single thirteen-row audit table.  The placeholder currently committed under
\texttt{results/claims/} is not evidence and cannot close this gate.

The campaign record consists of 515 checkpoint rows and associated raw or
derived per-chunk artifacts.  Part of that material was historically ignored
and must first be recovered read-only from the campaign machine, including
chunk~7 from checkpoint evidence although its journal write is absent.  The
release also requires byte-level equality of the scanner and independent CPU
ordered-prime streams for every chunk; matching congruence-class counts alone
does not establish identity.  The four exact historical scanner blobs are
already preserved in the publication repository with their revision mapping.

The separate registry \texttt{audit/rescan17.json} binds the preserved
post-campaign transcript for the first 17 chunks to its full Git commit, source
blob, raw SHA-256, and eight independent counts committed before the run
completed.  A focused public provenance record contains those preregistered
counts without importing the unrelated internal journal from which they were
extracted.  Running \texttt{tools/rescan17\_audit.py} parses the historical
format fail-closed and generates Table~\ref{tab:rescan17}.  Because the old run
did not preserve separate standard streams, a binary hash, or per-chunk rows,
the registry records those fields as unavailable rather than reconstructing
them.  This aggregate audit does not substitute for the 515-chunk
ordered-prime identity gate.

The cyclic $W_c$ pilot and its proof artifacts are outside the scientific
scope of this records paper and are intentionally absent from this repository.
The archived two-colour phase-2 source is credited and located, but a verbatim
third-party web capture is not redistributed unless its rights status is
resolved.  These provenance metadata support priority attribution; they do not
turn an archived claim of exhaustive coverage into a newly audited fact.

Before submission, a clean release candidate must rebuild the scanner and
verifiers, execute and validate the thirteen claim manifests, regenerate every
table, and compile this PDF.  After the pre-freeze gates pass and the DOI
metadata are final, the entire intended payload is staged.  Immediately before
the inner-freeze transaction, the staged release form of \texttt{STATUS.json}
sets \texttt{repository\_state} to \texttt{archived\_release}, sets
\texttt{citable\_release} true, and marks every gate passed.  For the DOI gate,
that transient staged state means only that the metadata are final and the
inner-freeze transaction is ready; it does not assert that a tag, archive,
deposit, or receipt is already public.

The transaction \texttt{tools/freeze\_release.py create} validates that staged
candidate, rebuilds the PDF, and writes the top-level
\texttt{MANIFEST.sha256}.  This is the first checksum layer and binds the
tracked payload files: the PDF, source files, certificates, logs, and data.  If
the transaction fails, the staged release state is neither committed nor
published.  Only after success is the manifest staged with the payload, which
is then committed, tagged locally, and rechecked.  A deterministic archive is
subsequently created from that checked immutable local tag and rechecked
byte-for-byte by \texttt{tools/build\_release\_archive.py}; this is the second
checksum layer.  The archive's own byte length and SHA-256, together with the
post-tag draft-deposit comparison, belong in an external deposit receipt or
release record.  They cannot be inserted into the archive they hash without a
cycle.  The fail-closed utility \texttt{tools/deposit\_receipt.py} creates and
rechecks that JSON only at an absolute path outside the repository, and the
receipt is not embedded evidence for the DOI gate.

Four fixed headline certificates must additionally be checked on independently
administered hardware with complete outputs preserved.  After the scientific,
replication, and rights gates pass, the DOI can be reserved and inserted into
the exact final candidate.  The successful inner-freeze transaction precedes
the commit and local tag; the checked tagged tree is then archived and compared
with the draft upload.  The tag and deposit are made public only if that
comparison succeeds.  Until then this manuscript does not claim that a
complete citable artifact is available.
